\documentclass[11pt]{article}
\usepackage[margin=0.72in]{geometry}
\usepackage{fontspec}
\usepackage{setspace}
\usepackage{hyperref}
\hypersetup{colorlinks=true, urlcolor=blue, linkcolor=black}
\pagestyle{empty}
\setlength{\parskip}{0.6em}
\setlength{\parindent}{0pt}
\setstretch{0.93}
\setmainfont{Times New Roman}

\begin{document}

\begin{flushleft}
Ke Cheng HSU (Clio) \\
Department of Economics, The Hong Kong University of Science and Technology \\
Clear Water Bay, Kowloon, Hong Kong $\cdot$ +852 6840 9753 \\
\href{mailto:khsu@connect.ust.hk}{khsu@connect.ust.hk} $\cdot$ \href{https://cliohsu1997.github.io/}{cliohsu1997.github.io}
\end{flushleft}

16 November 2026

\vspace{0.08em}
Human Resources Office \\
Hong Kong Chu Hai College \\
Hong Kong

\vspace{0.08em}
Dear Members of the Search Committee,

I am writing to apply for the Research Fellow / Research Associate position under Project Ref.\ No.\ CHC/RMGS/19/001 (Belt and Road Studies -- Development of Economy, Industries and Trade in Belt and Road Countries). I am a Ph.D.\ candidate in Economics at HKUST, expect to graduate in June 2027, and can start as early as February 2027. I am a Hong Kong permanent resident with the right to work in Hong Kong, a native Mandarin speaker, and professionally fluent in English. Since 2020 I have worked full-time on empirical research through my M.Phil.\ and Ph.D., including a research assistantship at National Taiwan University---experience that maps closely to assisting a project leader on data-intensive economic research, drafting, and research events.

What I offer this project is research capacity on economy and industry questions, bilingual communication, and reliable project support. My sole-authored job market paper, ``Designing Return Policy in E-Commerce: A Trade-off between Purchase and Return Efficiency,'' studies how return policy affects purchases, returns, and welfare using large Taobao apparel transaction data from China. I design the data pipeline, estimate a structural model, and quantify counterfactual policies: a menu with a moderate restocking fee raises social welfare by about 2 million RMB (roughly a 5\% gain) and seller profit by about 1.2 million RMB relative to status-quo insurance policies. With Pak Hung Au, I study product design under heterogeneous search costs. With Kohei Kawaguchi and Kei Kawai, I estimate pricing and inventory under demand uncertainty from Japanese bookstore point-of-sale records spanning tens of billions of observations. These projects show that I can collect and analyze complex industry data, compare policy scenarios with clear numbers, and write results for a professional audience---the core research tasks listed in the posting.

I am also prepared for the liaison and organization side of the role. I have presented at the Asia-Pacific Industrial Organization Conference, the Econometric Society Australasia Meeting, and the China Economics Annual Conference in 2026, so conference and seminar settings are familiar. As Teaching Assistant Coordinator in the Business School at HKUST, I coordinated a teaching team and held materials to one standard before release; as lecturer and TA my evaluations are 4.53/5 and 4.72/5. I can support stakeholder communication and research activities in Mandarin and English, and I am motivated to apply this toolkit to Belt and Road questions on economic development, industries, and trade under the project leader's direction.

My present monthly compensation is HKD\,20,000 (postgraduate studentship). My expected monthly salary is HKD\,25,000--30,000. I attach my CV. Thank you for your consideration. You can reach me at \href{mailto:khsu@connect.ust.hk}{khsu@connect.ust.hk} or +852 6840 9753.

\vspace{0.18em}
Sincerely,

\vspace{0.18em}
Ke Cheng HSU \\
Ph.D.\ Candidate, Department of Economics, HKUST

\end{document}
